\documentclass[10pt,a4paper]{amsart}
\usepackage[margin=19mm]{geometry}
\usepackage{amsmath,amssymb,mathtools}
\usepackage[scr=rsfs,cal=euler]{mathalfa}
\usepackage{xfrac}
\usepackage[unicode,hidelinks,bookmarks=false]{hyperref}
\newcommand{\ie}{i.e.\ }
\newcommand{\cf}{cf.\ }
\newcommand{\resp}{respectively\ }
\newcommand{\Subsection}{\subsection}
\providecommand{\nobreakdash}{\nobreak\hbox{-}}
\setlength{\emergencystretch}{2em}
\multlinegap0pt
\usepackage{bm}
\usepackage{bbm}
\usepackage{dsfont}
\usepackage{xspace}
\usepackage{cancel}
\usepackage{mathtools}
\usepackage{amsthm}
\makeatletter
\@ifundefined{theorem}{\newtheorem{theorem}{Theorem}}{}
\makeatother
\def\h{\mathbb H}
\def\l{\mathbb L}
\def\r{\mathbb R}
\def\s{\mathbb S}
\title[Classification of ruled surfaces in Lorentz-Minkowski space ]{Classification of ruled surfaces in Lorentz-Minkowski space that are stationary for the moment of inertia}
\author{Muhittin Evren Aydin, Rafael L\'opez}
\date{Source: arXiv:2508.17317v1 (CC BY 4.0)}
\begin{document}
\maketitle

\noindent\emph{Source and licence.} Classification of ruled surfaces in Lorentz-Minkowski space that are stationary for the moment of inertia. Version arXiv:2508.17317v1 (CC BY 4.0), distributed under the Creative Commons Attribution 4.0 International licence, as recorded in the arXiv licence metadata for this exact version and shown on the abstract page https://arxiv.org/abs/2508.17317v1. Full rights notice: licenses/arxiv-2508.17317-CC-BY-4.0.txt.\par\smallskip

\noindent\textit{Editorial scope.} Counted unit: the Theorem whose source label is \texttt{th-nli} in the article (source0.tex lines 416--427, source label \texttt{th-nli}), delivered together with the article's own complete original proof of it (source0.tex lines 428--601, one \texttt{proof} environment of 174 lines). No mathematics is written, repaired or re-derived: statement and proof are the article's own text line for line. Required context retained uncounted: rp (lines 336--338); ss3 (lines 378--380). Every definition, display and label of those lines is retained unchanged. Omitted from this package and not redistributed: 1--335, 339--377, 381--415, 602--809. Nothing inside a retained excerpt is omitted or altered: every retained body is delivered exactly as the article's lines are written, so no omission receipt of any kind accompanies this package. Citations: the retained excerpts carry no citation command at all, so no bibliography entry of the article is transcribed and no reference is redistributed. Reference adaptation: none was needed and none was made; every reference inside the delivered unit resolves inside it, so no unresolved reference or citation is produced. Printed slips kept literal: no printed word, symbol, hypothesis or number of the counted statement or of its proof was altered. Layout and class: the article's own class is \texttt{amsart} with packages including amssymb, amsthm, graphicx, color, changes, url; the wrapper is the lane's pinned \texttt{amsart} editorial wrapper at 10pt on A4, which restates the theorem-like environments the retained text needs and adds the article's own macro definitions that the retained text needs (\texttt{\textbackslash h}, \texttt{\textbackslash l}, \texttt{\textbackslash r}, \texttt{\textbackslash s}), with extra wrapper packages (bm, bbm, dsfont, xspace, cancel, mathtools). The delivered numbering is the unit's own. Assets: no retained span contains a figure environment, a graphics-inclusion command, a PDF-page inclusion, a table or a TikZ picture; no image file is copied into this package, the package declares no asset dependency and no image file is redistributed with it. Licence: version arXiv:2508.17317v1 of this article is distributed under the Creative Commons Attribution 4.0 International licence, as recorded in the arXiv licence metadata for this exact version and shown on the abstract page https://arxiv.org/abs/2508.17317v1. A full rights notice is delivered at licenses/arxiv-2508.17317-CC-BY-4.0.txt. Characters: every retained excerpt is pure ASCII, so the delivered engine, XeLaTeX with its default Latin Modern text font, reports no missing character.
\begin{equation}
X(s,t)=\gamma(s)+tw(s), \quad s\in I, t\in \r,\label{rp}
\end{equation}

\begin{equation}
H_1=\alpha\epsilon\varepsilon(EG-F^2)\frac{(X_s,X_t,X)}{\langle X,X\rangle}, \label{ss3}
\end{equation}

\begin{theorem}[Case $w$ and $w'$ not lightlike] \label{th-nli}
Let $\Sigma \subset \mathcal{C}^\epsilon$ be a   non-cylindrical ruled surface parametrized by \eqref{rp} where $\langle w,w\rangle=\delta$, $\langle w',w' \rangle=\mu$, and $\delta,\mu\in \{-1,1 \}$. If $\Sigma$ is an $\alpha$-stationary surface, then $\alpha = \epsilon \varepsilon$, $\langle \gamma,w\rangle =0$, and $w=w(s)$ is a geodesic of $\s_1^2(1)$. Moreover, up to a linear isometry and a dilation of $\l^3$, one of the following holds:
\begin{enumerate}
\item $\mu=1$, $w(s)=(\cos s,\sin s,0)$, $\epsilon= \varepsilon$, and either
\begin{enumerate}
\item[\textnormal{(1a)}] $\gamma(s)=(0,0,e^{\pm s})$, or
\item[\textnormal{(1b)}]  $\gamma(s)=(-\sin s,\cos s,c_1 \cosh s+c_2 \sinh s)$ with $c_1^2+c_2^2=1$.
\end{enumerate}
\item $\mu=-1$,  $\epsilon= -\varepsilon$, $\gamma(s)=(\sinh s, c_1\cos s+c_2\sin s, \cosh s)$, with $c_1^2+c_2^2=1$ and $w(s)=(\cosh s,0,\sinh s)$.
 
\end{enumerate}
\end{theorem}

\begin{proof}
Without loss of generality, we can choose $\gamma$ to be the striction curve, so that $ \langle \gamma',w'\rangle=0$. The coefficients of  the first fundamental form are
$$
E=\langle \gamma',\gamma' \rangle +\mu t^2, \quad F=\langle \gamma',w\rangle, \quad G=\delta.
$$ 
Notice that $w(s)\in \s_1^2(1)$ or $w(s)\in \h^2(1)$, according to the sign of $\delta$, and denote by $\kappa_g^w=(w'',w',w)$ the geodesic curvature of $w(s)$. We will consider the orhonormal basis $\mathcal{B}=\{w,w',w\times w' \}$, where $\langle w\times w', w\times w' \rangle =-\delta \mu$. In terms of $\mathcal{B}$, we have
\begin{align*}
\gamma'&=\delta F w-\delta \mu Q (w \times w'), \\
w''&=\delta \mu (-w+\kappa_g^w(w\times w')), \\
w\times w''&=\mu \kappa_g^w w',
\end{align*}
where $Q=(\gamma',w,w')$ is the parameter of distribution. Differentiating $\gamma'$ with respect to $s$, we have
$$
\gamma''=\delta F'w +\delta (F-Q\kappa_g^w)w'-\delta \mu Q' (w\times w').
$$
Set $\beta=-\varepsilon \epsilon \in \{-1,1\}.$ Hence, Eq. \eqref{ss3} becomes
$$
\frac{-Q(F+Q\kappa_g^w)-\delta Q' t +\delta \kappa_g^wt^2}{\mu (\delta t^2-Q^2)}=-\alpha\beta\frac{(\gamma'+tw',w,\gamma)}{\langle \gamma,\gamma\rangle+2t\langle \gamma,w\rangle+\delta t^2}.
$$ 
This equation can be written  of the form $$\sum_{n=0}^4A_n(s)t^n=0.$$ 
Thus all coefficients   $A_n(s)$ vanish identically. For $n=3,4$, a computation yields 
$$
A_3=Q'-\alpha\beta\delta\mu(w',w,\gamma), \quad A_4=\kappa_g^w.
$$
From $A_4=0$,  $w=w(s)$ is a geodesic of $\s_1^2(1)$ or of $ \h^2(1)$. Also, it follows that $w\times w'$ is constant along $\gamma$. Since the condition $Q=0$ yields that $\Sigma$ is a vector plane, we will assume that $Q\neq 0$.

We write $\gamma$ in terms of $\mathcal{B}$, obtaining 
$$
\gamma=cw+aw'+b(w\times w'),
$$
where $a$, $b$, and $c$ are smooth functions defined on $I$.  Differentiating and using $\langle \gamma',w'\rangle=0$, we find $c=-a'$ on $I$. Consequently, 
\begin{align*}
\gamma&=-a'w+aw'+b(w\times w'), \\
\gamma'&=-(a''+\delta\mu a)w+b'(w\times w'),
\end{align*}
where $-(a''+\delta\mu a)=\delta F$ and $Q=-\delta\mu b'$. Thus, from $A_3=0$, we conclude
\begin{equation}
 b''+\alpha\beta\delta\mu b =0. \label{ode1}
\end{equation} 
On the other hand, the remaining coefficients $A_n$ are  
\begin{align*}
A_0&=QF\langle \gamma,\gamma \rangle +\alpha \beta \mu Q^2(\gamma',w,\gamma), \\
A_1&=2FQ\langle \gamma,w \rangle +\delta Q' \langle \gamma,\gamma \rangle+\alpha \beta \mu Q^2 (w',w,\gamma), \\
A_2&=2\delta Q'\langle \gamma,w \rangle +\delta QF-\alpha \beta \delta \mu  (\gamma',w,\gamma).
\end{align*}
We separate two cases:

\begin{enumerate}
\item  Case $F=\langle \gamma',w \rangle =0$. From $A_0=0$, we have $(\gamma',w,\gamma)=0$. By considering in $A_2=0$, we obtain that either $Q'=0$ or $\langle \gamma,w \rangle =0$. If $Q'=0$, then from \eqref{ode1} we arrive at $b=0$, i.e., $\Sigma$ is a plane, which is not our case. Thus,   $\langle \gamma,w \rangle =0$, which yields   $\gamma=b (w\times w')$. By using $A_1=0$, we find $bb''+\alpha \beta\delta\mu b'^2=0$. Comparing Eq. \eqref{ode1}, we obtain $b^2-\delta\mu b'^2=0$. Here, it must hold $\delta\mu=1$; otherwise we would have $b=0$ or $Q=0$, a contradiction. Since $\langle w,w'\rangle =0$, we arrive at $\delta=\mu=1.$ In addition, due to $b'=\pm b$, using \eqref{ode1}, we conclude $\alpha=-\beta$. This implies that either $b(s)=c_0e^s$ or $b(s)=c_0e^{-s}$, where $c_0 \in \r$ and $c_0 \neq 0$. Up to a dilation, we may assume that $c_0=1$. On the other hand, $EG-F^2=t^2-Q^2$ and $\langle X,X\rangle=\langle \gamma,\gamma\rangle+t^2=-b^2+t^2$, which implies $\epsilon=\varepsilon$. Therefore, after solving $w''+w=0$ and applying a linear isometry of $\l^3$, we find $w(s)=(\cos s, \sin s, 0)$, which completes the proof of the item (1a) of the theorem.

\item Case $F=\langle \gamma',w \rangle \neq 0$. If $Q'=0$, then from $A_3=0$ it follows that $(w',w,\gamma)=0$, that is, $b=0.$ This contradicts with $Q\neq 0$. From $A_0=0$ and $A_2=0$, we have $\langle \gamma, \gamma \rangle=-\alpha\beta\mu\frac{Q}{F}(\gamma',w,\gamma)$ and 
\begin{equation*}
\alpha\beta\mu(\gamma',w,\gamma)=2Q'\langle \gamma, w \rangle+QF,
\end{equation*}
respectively. This yields   
\begin{equation}
\langle \gamma, \gamma \rangle=-\frac{Q}{F}(2Q'\langle \gamma, w \rangle+QF). \label{qf}
\end{equation}
Inserting this in $A_1=0$, we obtain 
$$\langle \gamma, w \rangle (F^2-\delta Q'^2)=0.$$ 
We investigate the resulting cases separately.

\begin{enumerate}
\item Case $\langle \gamma, w \rangle = 0$. It follows that $a(s)=a_0\not=0$ is constant, and that $F = -\mu a_0$. Up to a dilation, we may assume that $a_0=1$. From the equations $A_0 = 0$ and $A_1=0$, we obtain $1= \delta b^2 + \alpha \beta \mu b'^2$  and $1= \delta b^2 - \mu b'^2$, respectively. Both equations imply  $\alpha \beta = -1$. Hence, there are two possibilities: $(\delta,\mu)=(1,1)$ and $(\delta,\mu)=(1,-1)$. If $\mu=1$, then, similarly to Case 1, we arrive at item (1b) of the theorem. Otherwise, $\mu=-1$,  solving $w''+\mu w=0$ and applying a linear isometry of $\l^3$ we obtain $w=(\cosh s, 0 , \sinh s).$ Consequently, after solving Eq.  \eqref{ode1}, we establish item (2) of the theorem. Similar to Case 1, it is direct to see that $EG-F^2=-\langle X,X\rangle$, i.e. $\epsilon=-\varepsilon$.

\item Case  $\langle \gamma, w \rangle \neq 0$. It follows $F^2-\delta Q'^2=0$, which requires $\delta=1$, that is, $F=\pm Q'$. Without loss of generality, suppose $F=Q'$, that is
\begin{equation}
a''+\mu a=\mu b''. \label{amub}
\end{equation}
On the other hand, from $A_2=0$, we conclude
\begin{equation}
2a'b+bb'-ab'=0. \label{ab}
\end{equation}
We will show that this case leads to a contradiction. The argument will be divided into cases according to the values of $\mu$ and $\alpha\beta$. Although the work is a bit tedious, the reasoning in each case is identical, where first we   solve the ODEs \eqref{ode1} and \eqref{amub}, and then substitute the solutions into \eqref{ab} obtaining an equation of functions on $s$ which are linearly independent. In particular, all coefficients of these functions must vanish identically. An analysis of theses coefficients gives the desired contradiction. Set $k=\alpha\beta$. To simplify the notation, let $c_i$ ($1 \leq i \leq 4$) denote the constants of integration that will appear below.
\begin{enumerate}
\item Case $\mu=1$ and $k=1$. The solutions of \eqref{ode1} and \eqref{amub} are
\begin{equation*}
\begin{split}
b(s)&=c_1\cos s+c_2 \sin s,\\
a(s)&=\frac12 ((-c_1+2c_3+c_2s)\cos s+(-c_1s+2c_4)\sin s).
\end{split}
\end{equation*}
Substituting   in \eqref{ab} yields a trigonometric equation whose coefficients must vanish. The coefficient of the term $s \cos^2 s$ is $-c_1^2-\frac12 c^2$, which is nonzero, leading to a contradiction.

\item Case $\mu=1$, $k >0$, and $k\neq1$. Now we have
\begin{equation*}
\begin{split}
 b(s)&=c_1\cos (\sqrt{k}s)+c_2 \sin (\sqrt{k}s),\\
a(s)&=\frac{k }{k-1}(c_1 \cos (\sqrt{k} s)+c_2 \sin (\sqrt{k} s))+c_3 \cos s+c_4 \sin s.
\end{split}
\end{equation*}
By substituting   in \eqref{ab}, we arrive at
$$
\begin{aligned}
0=& \sqrt{k}(-c_1^2+c_2^2)\sin(2\sqrt{k}s)
+ 2\sqrt{k}c_1c_2\cos(2\sqrt{k}s) \\
&+ (\sqrt{k}-1)(c_1c_3-c_2c_4)\sin((\sqrt{k}+1)s)
+ (\sqrt{k}+1)(c_1c_3+c_2c_4)\sin((\sqrt{k}-1)s) \\
&+ (\sqrt{k}+1)(c_1c_4-c_2c_3)\cos((\sqrt{k}-1)s)
- (\sqrt{k}-1)(c_1c_4+c_2c_3)\cos((\sqrt{k}+1)s).
\end{aligned}
$$
From the coefficients of the terms of $\sin(2\sqrt{k}s)$ and $\cos(2\sqrt{k}s)$, we find $c_1=c_2=0$, which is not possible.

\item Case $\mu=1$, $k <0$. We have 
\begin{equation*}
\begin{split}
b(s)&=c_1e^{\sqrt{-k}s}+c_2 e^{-\sqrt{-k}s},\\ 
a(s)&=\frac{k }{k-1}(c_1 e^{\sqrt{-k} s}+c_2 e^{-\sqrt{-k} s})+c_3 \cos (s)+c_4 \sin (s).
\end{split}
\end{equation*}
Using these expressions in \eqref{ab}, we obtain 
$$
\begin{aligned}
0=&
\sqrt{-k}\frac{2k-1}{k-1}( c_1^{2} e^{2\sqrt{-k}s}-c_2^{2} e^{-2\sqrt{-k}s}) \\
&+(c_1c_3(-2\sin s - \sqrt{-k}\cos s)+c_1c_4(2\cos s - \sqrt{-k}\sin s)) e^{\sqrt{-k}s} \\
&+(c_2c_3(-2\sin s +\sqrt{-k}\cos s)+c_2c_4(2\cos s + \sqrt{-k}\sin s))e^{-\sqrt{-k}s}.
\end{aligned}
$$
This leads to a contradiction, because the coefficients of the terms $e^{2\sqrt{-k}s}$ and $e^{-2\sqrt{-k}s}$ cannot both vanish simultaneously.

\item Case $\mu=-1$ and $k=1$. We have 
\begin{equation*}
\begin{split}
b(s)&=c_1e^s+c_2 e^{-s},\\
a(s) &= c_3 e^s + c_4  e^{-s} - \frac{c_1}{2} s e^s + \frac{c_2}{2} s e^{-s}.
\end{split}
\end{equation*}
Inserting these expressions in \eqref{ab}, we obtain
$$
(c_1 c_3 - \frac{c_1^2}{2}  s ) e^{2s}
+( 3c_2 c_3 - 3c_1 c_4 - 3c_1 c_2 s )
+ ( -c_2 c_4 - \frac{c_2^2}{2}  s ) e^{-2s}=0,
$$
which gives the contradiction $c_1=c_2=0$.

\item Case $\mu=-1$ and $k>0$, $k\neq 1$. We have
\begin{equation*}
\begin{split}
b(s)&=c_1e^{\sqrt{k}s}+c_2 e^{-{\sqrt{k}s}},\\
a(s)&= c_3e^{s} + c_4e^{-s}- \frac{k}{k-1}( c_1e^{\sqrt{k}s} + c_2e^{-\sqrt{k}s} ).
\end{split}
\end{equation*}
Considering these in \eqref{ab}, we obtain
$$
\begin{aligned}
0 &= (2 - \sqrt{k})\,c_1 c_3\, e^{(1+\sqrt{k})s}
+ (2 + \sqrt{k})\,c_2 c_3\, e^{(1-\sqrt{k})s} \\
&\quad - (2 + \sqrt{k})\,c_1 c_4\, e^{(-1+\sqrt{k})s}
- (2 - \sqrt{k})\,c_2 c_4\, e^{-(1+\sqrt{k})s} \\
&\quad - \frac{\sqrt{k}}{k-1} c_1^2 e^{2\sqrt{k} s}
+ \frac{\sqrt{k}}{k-1} c_2^2 e^{-2\sqrt{k} s},
\end{aligned}
$$
where the coefficients must vanish. But this gives  $c_i = 0$, a contradiction.

\item Case $\mu=-1$ and $k<0$. We have 
\begin{equation*}
\begin{split}
b(s)&=c_1\cos (\sqrt{-k}s)+c_2 \sin(\sqrt{-k}s),\\
a(s) &= c_3 e^{s} + c_4 e^{-s}+ \frac{k}{k-1} ( c_1 \cos( \sqrt{-k} s )+c_2 \sin( \sqrt{-k}s ) .
\end{split}
\end{equation*}
Substituting these into \eqref{ab} and focusing only on the purely trigonometric terms, we find  
$$
\frac{\sqrt{-k}(1-2k)}{k-1} ( \frac12(c_{1}^2 - c_{2}^2) \sin(2 \sqrt{-k}s) - c_{1} c_{2} \cos(2 \sqrt{-k} s)).
$$
This gives $c_1=c_2=0$,  a contradiction.

\end{enumerate}
\end{enumerate}
\end{enumerate}
\end{proof}

\end{document}
